\chapter{Introduction}
\label{ch:introduction}

\section{Background}
\label{sec:background}

Virtual try-on (VTON) is a computer vision and generative AI application that allows users to visualize how a garment would look on their body without physically wearing it. The technology has gained significant commercial traction in the fashion retail industry, where \textit{smart mirrors}---large interactive displays equipped with cameras---enable shoppers to preview outfits in stores, malls, and trade exhibitions.

Traditional fitting rooms create bottlenecks during peak shopping hours, increase inventory handling costs, and limit the number of garments a customer can try in a single visit. Virtual try-on addresses these limitations by digitizing the fitting experience. However, deploying state-of-the-art VTON in production environments presents substantial engineering challenges:

\begin{itemize}[leftmargin=*]
  \item Neural inference on shared GPU clouds is \textbf{slow and queue-dependent}, with unpredictable wall-clock latency.
  \item Model outputs are \textbf{not guaranteed} to be valid try-on composites; models may return unchanged person images, raw garment images, or corrupted data.
  \item Kiosk software must handle \textbf{many sequential users} without cross-session UI corruption from stale asynchronous jobs.
  \item Exhibition and retail deployments require \textbf{continuous uptime} where downtime is unacceptable.
\end{itemize}

\section{Project Overview}
\label{sec:project-overview}

The \textbf{Smart Mirror AI Virtual Try-On System} is a full-stack web application designed to address these challenges. The system enables users to browse a digital garment catalog, capture a live webcam photograph, and receive a composite image approximating how the selected garment would appear on their body.

The project comprises three major components:

\begin{enumerate}[leftmargin=*]
  \item \textbf{Frontend (React Kiosk):} A single-page application built with React 18, Vite 5, and Tailwind CSS 3, providing a guided kiosk flow with catalog browsing, camera capture, result display, and an optional AI stylist assistant.
  \item \textbf{Backend (Node.js/Express):} An Express 5 server orchestrating catalog resolution, image preprocessing, Replicate cloud predictions, local compositing fallbacks, and output validation.
  \item \textbf{Data Layer:} MongoDB for garment metadata storage and static PNG assets served from \texttt{frontend/public/dataset/}.
\end{enumerate}

The AI pipeline uses cloud-based diffusion try-on models (CATVTON-Flux and IDM-VTON via Replicate) as the primary inference path, supplemented by a multi-layer deterministic fallback chain for reliability.

\section{Motivation}
\label{sec:motivation}

The motivation for this project stems from the gap between research-grade virtual try-on systems and deployable retail solutions. While academic VTON models achieve impressive visual quality, they are typically evaluated in controlled laboratory settings with dedicated GPU hardware. Real-world kiosk deployments face network variability, shared cloud GPU queues, and non-technical end users who expect immediate, reliable results.

This project aims to bridge that gap by engineering a \textbf{reliability envelope} around cloud AI inference---combining generative quality with deterministic fallbacks, instant previews, session-safe async UI behavior, and strict output validation.

\section{Scope of the Project}
\label{sec:scope}

The scope of this project includes:

\begin{itemize}[leftmargin=*]
  \item Design and implementation of a kiosk-style virtual try-on web application.
  \item Integration with Replicate-hosted CATVTON-Flux and IDM-VTON models.
  \item Development of a multi-layer exhibition fallback chain (ComfyUI, MoveNet pose overlay, Sharp compositing).
  \item REST API design for catalog, try-on, and AI stylist services.
  \item Session isolation and async prediction polling on the frontend.
  \item MongoDB-based catalog management with offline mock fallback.
\end{itemize}

\textbf{Out of scope:} 3D body mesh reconstruction, physics-based cloth simulation, mobile application development, payment integration, and user account management.

\section{Objectives}
\label{sec:objectives}

\subsection{Functional Objectives}

\begin{table}[H]
  \centering
  \caption{Functional objectives of the Smart Mirror system}
  \label{tab:functional-objectives}
  \begin{tabular}{@{}clp{9cm}@{}}
    \toprule
    \textbf{ID} & \textbf{Objective} & \textbf{Description} \\
    \midrule
    F1 & Catalog browsing & Browse garments by gender and category with images from catalog API \\
    F2 & Webcam capture & Capture live webcam frame and send to server as normalized data URL \\
    F3 & Virtual try-on & Run virtual try-on using cloud AI and/or local fallbacks \\
    F4 & Instant preview & Show immediate Sharp composite preview while async AI processes \\
    F5 & Async polling & Poll async Replicate prediction until success or failure \\
    F6 & Result display & Display one final validated processed image on result screen \\
    F7 & AI stylist & Optional outfit suggestions via Ollama-powered stylist assistant \\
    \bottomrule
  \end{tabular}
\end{table}

\subsection{Non-Functional Objectives}

\begin{table}[H]
  \centering
  \caption{Non-functional objectives of the Smart Mirror system}
  \label{tab:nonfunctional-objectives}
  \begin{tabular}{@{}clp{9cm}@{}}
    \toprule
    \textbf{ID} & \textbf{Objective} & \textbf{Description} \\
    \midrule
    NF1 & Reliability & Bounded worst-case output via fallback chain and validation \\
    NF2 & Session safety & Stale async jobs must not overwrite new user sessions \\
    NF3 & Operability & Environment-driven model selection (CATVTON vs IDM-VTON) \\
    NF4 & Security & API tokens and HuggingFace tokens kept server-side only \\
    NF5 & Demonstrability & Exhibition mode hides failure-heavy UX for public demos \\
    \bottomrule
  \end{tabular}
\end{table}

\section{Problem Statement}
\label{sec:problem-statement}

Design and implement a smart mirror system that:

\begin{enumerate}[leftmargin=*]
  \item Accepts a live person image and a catalog garment image as inputs.
  \item Produces a plausible try-on result within acceptable demonstration latency.
  \item Never breaks the kiosk user interface with null images, wrong images, or session bleed from prior users.
  \item Continues operating when cloud AI is slow or fails, using deterministic local fallbacks.
\end{enumerate}

\section{Intended Users and Deployment Context}
\label{sec:users}

\begin{table}[H]
  \centering
  \caption{Deployment contexts for the Smart Mirror system}
  \label{tab:deployment-context}
  \begin{tabular}{@{}lp{10cm}@{}}
    \toprule
    \textbf{Context} & \textbf{Usage} \\
    \midrule
    Retail flagship store & Shoppers try outfits on a floor-standing smart mirror \\
    Trade show / exhibition & Short sessions with network and GPU variability \\
    College / lab demo & Development with mock catalog when MongoDB is offline \\
    \bottomrule
  \end{tabular}
\end{table}

\section{Organization of the Report}
\label{sec:organization}

This report is organized as follows:

\begin{itemize}[leftmargin=*]
  \item \textbf{Chapter 2} presents the background, literature survey, and technology choices.
  \item \textbf{Chapter 3} describes the Software Requirements Specification (SRS).
  \item \textbf{Chapter 4} covers system modeling, architecture, and implementation details.
  \item \textbf{Chapter 5} discusses testing methodology, test cases, and validation results.
  \item \textbf{Chapter 6} provides conclusions and directions for future work.
  \item \textbf{Appendix} contains environment variables, API reference, and module dependency graphs.
\end{itemize}
